\chapter{Injury Management}

\section{Burns and Scalds}

Burns are categorized by depth of tissue damage: first-degree (superficial), second-degree (partial thickness), and third-degree (full thickness).

\subsection{Burn Classification and Signs}

\begin{center}
\small
\begin{tabularx}{\textwidth}{@{}>{\raggedright\arraybackslash}l>{\raggedright\arraybackslash}X>{\raggedright\arraybackslash}X>{\raggedright\arraybackslash}X>{\raggedright\arraybackslash}X@{}}

\hline
\textbf{Degree} & \textbf{Tissue Involved} & \textbf{Appearance} & \textbf{Pain Level} & \textbf{Healing Time} \\
\hline
First & Epidermis only & Red, dry, painful & High & 3-7 days \\
Second & Dermis also & Blisters, red/swollen, very painful & High & 2-3 weeks \\
Third & full thickness (all layers) & white, blackened, charred or leathery, may be painless & Low (nerves destroyed) & Months to years (requires grafting) \\
Fourth & deeper structures (bone/muscle) & Charring, exposed deeper tissues & Variable & Life-threatening \\
\hline
\end{tabularx}
\end{center}

\textbf{WARNING:} \textbf{Do not} apply ice, butter, oils, greases, or toothpaste to burns — these can cause further tissue damage and increase infection risk. Do NOT break blisters.

\subsection{Burn Extent: The Rule of Nines}

In adults, the Rule of Nines estimates how much body surface area a burn covers, which guides how serious the burn is and whether hospital care is needed:

\begin{center}
\begin{tabular}{@{}ll@{}}
\hline
\textbf{Body Part} & \textbf{Body Surface Area} \\
\hline
Head and neck & 9\% \\
Each arm (front and back) & 9\% \\
Front of trunk (chest and abdomen) & 18\% \\
Back of trunk (back and buttocks) & 18\% \\
Each leg (front and back) & 18\% \\
Genital area & 1\% \\
\hline
\end{tabular}
\end{center}

The Rule of Nines applies to adults. In children, the head is proportionally larger and the legs smaller, so the percentages differ — for a child, the head is about 14\% and each leg about 14\%. The rule is an approximation for quick field estimation only; healthcare providers use more precise methods. Burns over roughly 10\% or more of body surface area in an adult (or a smaller area in a child) generally warrant emergency medical care.

\subsection{First Aid for Minor Burns (First-Degree and Small Second-Degree)}

\begin{enumerate}[leftmargin=*]
\item Remove the person from the source of heat
\item Run cool (not cold) water over the burn for 10-20 minutes OR cover with a clean, wet cloth if running water is unavailable \cite{aha_arc_2024}
\item Gently remove constrictive items (jewelry, watches, tight clothing) before swelling occurs
\end{enumerate}
\quicknote{\textbf{Timing}: Cooling should begin as soon as possible after the burn; effectiveness diminishes after 30 minutes.}

\begin{enumerate}[leftmargin=*]
\item Cover the burn with a sterile, non-adhesive dressing or clean cloth
\item Give pain relief (acetaminophen or ibuprofen if appropriate and safe)
\end{enumerate}
\subsubsection{Chemical Burns}

\begin{itemize}[leftmargin=*]
  \item Remove contaminated clothing carefully (cut away, do not pull off if stuck to skin)
  \item Flush copiously with running water for at least 20 minutes (more for certain chemicals like lime powder or metallic sodium which react violently with water -- brush off first, then flush)
  \item Identify the chemical and inform medical personnel
\end{itemize}

\subsubsection{Electrical Burns}

\begin{itemize}[leftmargin=*]
  \item Ensure power source is OFF before touching the victim
  \item Use a non-conductive object to move the victim away from the source if power cannot be turned off
  \item Look for entry and exit burns (often more serious than they appear internally)
  \item Assume potential cardiac effects even if minor external burns -- seek medical evaluation
\end{itemize}

\subsection{Major Burns Requiring Emergency Care}

Seek immediate medical attention if:

\begin{itemize}[leftmargin=*]
  \item Burn covers an area larger than the person's palm (adult) or extends across a joint
  \item Burn is third-degree (white, blackened, leather-like) in appearance
  \item Burn involves face, hands, feet, genitals, or major joints
  \item Burn is caused by electricity or chemicals
  \item The person shows signs of shock or distress
\end{itemize}

\section{Fractures and Bone Injuries}

A fracture is a break in the bone. An open (compound) fracture involves broken bone piercing through the skin, carrying high risk of infection.

\subsection{Signs and Symptoms}

\begin{itemize}[leftmargin=*]
  \item Pain, especially when moving or touching the area
  \item Swelling and bruising
  \item Deformity (arm/leg looks crooked or unnatural)
  \item Inability to bear weight or use the affected part
  \item A grinding sensation or sound
  \item In open fractures: visible bone ends or bleeding
\end{itemize}

\textbf{WARNING:} \textbf{Do not attempt to realign or straighten a fracture}. Moving the broken bone can cause further injury to blood vessels, nerves, and soft tissue.

\subsection{Management}

\begin{enumerate}[leftmargin=*]
\item Keep the injured person still and calm
\item Control any bleeding with direct pressure using a dressing (for open fractures)
\item Immobilize the injured part using a splint. Splint the joints above and below the fracture if possible
\end{enumerate}
\textbf{WARNING:} \textbf{CAUTION:} Splint Material: Use rigid materials such as board, rolled-up magazine, cardboard, or commercial splint. Pad between splint and skin to prevent pressure sores. Check circulation distal to the splint (color, warmth, sensation). Loosen if excessive tightness develops.

\begin{enumerate}[leftmargin=*]
\item Apply cold packs (wrapped in cloth) to reduce swelling and pain -- avoid direct contact with skin
\item Seek emergency medical transportation for significant fractures
\end{enumerate}

\subsection{Open (Compound) Fracture Protocol}

An open fracture has a high risk of infection and must be treated with extra care:

\begin{enumerate}[leftmargin=*]
\item Treat the wound as an emergency: call emergency services immediately
\item Do NOT push protruding bone back in or attempt to realign the limb
\item Cover the wound with a clean, sterile (or at least clean) dressing -- use a moist sterile dressing if available
\item Control any bleeding around the wound with gentle direct pressure on the surrounding skin, not on the protruding bone
\item Splint the limb in the position found, immobilizing the joints above and below, without repositioning the bone
\item Keep the person still and warm, and monitor for shock -- open fractures can cause significant blood loss
\item Tell responders the time of injury and what was done
\end{enumerate}

\textbf{WARNING:} \textbf{Do not probe, clean inside, or irrigate an open fracture wound. Leave the exposed bone and deep tissues to professional medical care.}

\section{Sprains and Strains}

\begin{description}
  \item[Sprain:] Stretching or tearing of ligaments (tissue connecting bones at a joint)
  \item[Strain:] Stretching or tearing of muscles or tendons (tissue connecting muscle to bone)
\end{description}

\subsection{The R.I.C.E. Protocol}

\begin{itemize}[leftmargin=*]
  \item \textbf{R}est: Stop activity immediately, avoid putting weight on the injured area
  \item \textbf{I}ce: Apply cold pack for 15-20 minutes every 2-3 hours during the first 48 hours
  \item \textbf{C}ompression: Wrap with elastic bandage to reduce swelling (snug but not tight enough to cut off circulation)
  \item \textbf{E}levation: Raise the injured area above heart level when possible
\end{itemize}

\quicknote{\textbf{When to Seek Medical Attention}: If pain persists beyond a few days, there is significant deformity, you cannot bear weight, or numbness develops, get evaluated for possible fracture or ligament tear.}

\section{Head Injuries}

Any blow to the head that causes loss of consciousness or other symptoms should be taken seriously.

\quicknote{\textbf{Fall Mechanism}: For falls from stairs or significant height (> standing height), maintain higher suspicion for serious head, neck, or spinal injury. Always assume potential spinal injury until cleared by medical evaluation when the fall involved multiple steps, hard surface landing, or loss of consciousness.}

\subsection{Concussion Signs and Symptoms}

\begin{itemize}[leftmargin=*]
  \item Headache or pressure in head
  \item Nausea or vomiting
  \item Dizziness or balance problems
  \item Blurred or double vision
  \item Sensitivity to light or noise
  \item Feeling slowed down or "not right"
  \item Difficulty concentrating or remembering
  \item Confusion, disorientation
  \item Brief loss of consciousness (any LOC warrants medical evaluation)
\end{itemize}

\textbf{WARNING}: \textbf{Watch for Delayed Worsening}: Some intracranial injuries develop symptoms hours later. Keep under observation for at least 24 hours after significant head injury.

\subsection{Management}

\begin{enumerate}[leftmargin=*]
  \item Assess responsiveness and breathing
  \item If unresponsive and not breathing, begin CPR as indicated
  \item If responsive, keep the person still and stabilize their neck manually until spinal injury is ruled out
  \item Do NOT give food, drink, or medication (risk of surgery needing fasting status)
  \item Monitor closely for changes in consciousness, repeated vomiting, worsening headache, unequal pupils, clear fluid draining from ears/nose, or seizures
\end{enumerate}
\textbf{WARNING}: \textbf{Return-to-Play/Activity}: After a diagnosed concussion, gradual return to normal activities only after symptom-free at rest and then gradually with increasing exertion, under medical guidance. Never return on same day of injury.

\subsection{Facial and Oral Trauma}

Falls and accidents can cause bleeding from the mouth, nose, or face. Most oral bleeding is from soft tissue injuries (lips, tongue, gums) and can be controlled with direct pressure. However, persistent or heavy bleeding may indicate a fractured jaw, severe laceration, or deeper injury.

\subsubsection{Controlling Oral Bleeding}

\begin{enumerate}[leftmargin=*]
  \item Have the person sit up and lean slightly forward (to prevent blood flowing down the throat)
  \item Apply direct pressure to the bleeding site using a clean cloth or sterile gauze — have the person bite gently on the pressure pack if the injury is inside the mouth
  \item If bleeding persists, apply an ice pack wrapped in cloth to the outside of the cheek over the bleeding area
  \item Do NOT try to push a bleeding tongue back into the mouth — this can obstruct the airway
  \item If possible, have the person spit into a basin rather than swallow blood (swallowing can cause nausea/vomiting)
\end{enumerate}
\subsubsection{When to Suspect Fracture}

Seek emergency care if any of the following are present:
\begin{itemize}[leftmargin=*]
  \item Inability to close the mouth normally (malocclusion — upper and lower teeth don't fit together properly)
  \item Crepitus (grating sensation) when moving the jaw
  \item Significant swelling or bruising around the jaw, eyes, or cheeks within hours
  \item Numbness of the lower lip or chin (possible nerve involvement)
  \item Persistent heavy bleeding despite 15 minutes of direct pressure
  \item Blood draining freely from the ear or clear fluid from the nose (possible basal skull fracture)
\end{itemize}

\subsubsection{Knocked-Out Tooth (Avulsed Tooth)}

If a permanent tooth is completely knocked out:
\begin{enumerate}[leftmargin=*]
  \item Pick up the tooth by the crown (white part) — DO NOT touch or scrub the root
  \item If dirty, rinse briefly in water only (no soap, no scrubbing)
  \item Attempt to reinsert the tooth into the socket if possible and safe — have the person bite gently on gauze to hold it
  \item If reinsertion is not possible, store the tooth in milk, saline solution, or saliva (place under the person's tongue). Do NOT store in dry conditions or plain water
  \item Seek dental care immediately — best outcomes within 30 minutes to 1 hour
\end{enumerate}
For children with knocked-out baby teeth, do NOT attempt reinsertion — seek dental evaluation instead.



\subsection{Nosebleeds (Epistaxis)}

Nosebleeds are common after falls, head trauma, or facial injuries. Most anterior nosebleeds can be controlled with simple measures.

\subsubsection{Immediate Steps}

\begin{enumerate}[leftmargin=*]
  \item Have the person sit up and lean slightly forward — this prevents blood from flowing down the throat (which causes nausea/vomiting)
  \item Pinch the soft part of the nose (just below the bony ridge) firmly between the thumb and index finger
  \item Maintain continuous pressure for at least 10–15 minutes without releasing to check
  \item Apply an ice pack wrapped in cloth to the bridge of the nose or back of the neck (constriction may help)
  \item After bleeding stops, advise avoiding nose-blowing, bending over, or strenuous activity for several hours
\end{enumerate}
\subsubsection{When to Seek Emergency Care}

Call emergency services if:
\begin{itemize}[leftmargin=*]
  \item Bleeding continues despite 15–20 minutes of firm pressure
  \item Bleeding is heavy and flows freely down the back of the throat
  \item The nosebleed follows a significant head injury or facial trauma
  \item The person has a known bleeding disorder or is on blood-thinning medication
  \item Repeated nosebleeds occur without an obvious cause
\end{itemize}

\noindent\textbf{Note}: Do NOT stuff tissue or gauze deeply into the nostril unless trained to do so — improper placement can worsen bleeding or cause infection. If you must pack loosely, seek medical help promptly.

\section{Eye Injuries}

\subsection{Foreign Body in Eye}

\begin{itemize}[leftmargin=*]
  \item Have the person blink several times to encourage natural tear flushing
  \item Pull lower eyelid down and upper lid up to check for objects
  \item If visible on the white of eye or inside of lid, gently remove with moistened cotton swab or corner of clean tissue
  \item Do NOT rub the eye
  \item If object is embedded or under the eyelid AND cannot be easily removed, seek medical help immediately
  \item Cover both eyes (movement of one eye can move the other) during transport
\end{itemize}

\subsection{Chemical Eye Injury}

\begin{enumerate}[leftmargin=*]
\item Immediately flush eye with clean water or saline for at least 15-20 minutes
\item Hold eyelids open during flushing to ensure thorough irrigation
\item Flush from nose outward (to avoid contaminating other eye)
\item Seek emergency medical care even after thorough flushing
\end{enumerate}
\textbf{WARNING:} \textbf{Alkali substances} (like ammonia, lye, bleach) require prolonged flushing as they can continue penetrating tissue. Do not try to neutralize chemically.

\subsection{Emergency Eye Irrigation Technique}

When flushing a contaminated eye, technique matters more than speed alone:

\begin{enumerate}[leftmargin=*]
\item Remove contact lenses immediately if they are in the eye (they can trap the chemical against the eye); do not delay flushing to remove them
\item Tilt the head so the affected eye is down and the water runs from the inner corner (nose side) toward the outer corner, keeping the chemical away from the other eye
\item Hold the eyelids open with the fingers of one hand (forefinger on upper lid, thumb on lower lid) while pouring with the other, or use a clean container of water
\item If no running water or eyewash station is available, pour clean water from a bottle, cup, or faucet over the eye for the full duration
\item For chemicals, flush continuously for at least 15-20 minutes; for alkali (lye, cement, drain cleaner), flush longer -- 30 minutes or more -- because these continue to burn deep tissue
\item Turn the person's head side to side and use a container with a steady stream so the water rinses the entire eye surface
\item Do NOT bandage the eye tightly or let the person rub it; seek emergency medical care after flushing
\end{enumerate}

\textbf{WARNING:} \textbf{Do not use any acid or alkali to "neutralize" a chemical in the eye.} Water or saline only.

\section{Spinal Injuries}

Suspect spinal injury with: mechanism of injury (high fall, diving accident, car crash, motor vehicle-pedestrian collision), any neurological symptoms (numbness, weakness, tingling), pain along spine, or intoxication preventing proper assessment.

\textbf{WARNING:} \textbf{Do not move the person unless absolutely necessary} (immediate danger like fire or water). Movement could worsen an existing spinal injury leading to permanent paralysis.

\subsection{Management}

\begin{itemize}[leftmargin=*]
  \item Manually stabilize the head and neck in neutral position (do not tilt, rotate, or extend the neck)
  \item Ask anyone nearby to maintain manual stabilization while you call for help
  \item Only move if the scene becomes unsafe (use log-roll technique with assistance if absolutely required)
  \item Place rolled towels or clothing on both sides of the head to prevent movement during transport
  \item Monitor breathing and consciousness continuously
\end{itemize}

\section{Chest Injuries}

\subsection{Sucking Chest Wound (Open Pneumothorax)}

A wound to the chest wall that lets air enter the space around the lung can cause the lung to collapse. It may produce a sucking sound with each breath. These wounds are life-threatening and need immediate emergency care.

\begin{enumerate}[leftmargin=*]
\item Call emergency services immediately
\item Keep the person as still as possible in a position of comfort, usually sitting up or half-sitting, leaning toward the injured side
\item If the wound is open and air can pass through it, cover it with a gloved hand immediately while you prepare a dressing
\item Apply an occlusive dressing: use a plastic sheet, plastic wrap, or the sterile packaging from a dressing. Tape it down on THREE sides only, leaving one corner untaped
\item If the person becomes worse (increasing trouble breathing), lift one corner of the dressing briefly to release trapped air, then re-seal
\item Monitor breathing closely; if the person becomes unresponsive and not breathing normally, begin CPR
\end{enumerate}

\textbf{WARNING:} \textbf{The three-sided dressing creates a "flutter valve":} it seals the wound when the person breathes in, keeping air out, but lifts to let trapped air escape when they breathe out. Taping all four sides can trap air and collapse the lung further (tension pneumothorax).

\textbf{WARNING:} \textbf{Do not remove an object embedded in the chest.} Stabilize it in place with bulky dressings exactly as described in Chapter~\ref{ch:bleeding}, Section~\ref{sec:embedded_objects}.

\subsection{Rib and Flail Chest}

\begin{itemize}[leftmargin=*]
  \item Suspect rib fracture after chest trauma with pain on breathing or pressing the area
  \item Have the person hold a pillow or folded clothing against the injured side and support the arm over it while awaiting help -- do NOT tape or tightly bind the chest
  \item In a flail chest (two or more ribs broken in two places, causing a segment that moves opposite to the chest), suspect significant internal injury -- call emergency services and keep the person still
\end{itemize}

\section{Soft Tissue Injuries (Contusions/Abrasions)}

\subsection{Bruises (Contusions)}

Manage with the R.I.C.E. protocol (Rest, Ice, Compression, Elevation). Most bruises resolve within 1-2 weeks. Seek medical attention if: large or expanding bruise without apparent cause, bruise in sensitive areas (eye socket, abdomen, groin), bruise accompanied by severe pain, or inability to move nearby joint normally.

\subsection{Scrapes (Abrasions)}

\begin{itemize}[leftmargin=*]
  \item Wash gently around wound with mild soap and water (avoid getting soap directly in wound as it irritates)
  \item Remove surface debris with tweezers (cleaned with alcohol) or rinse away loose particles
  \item Apply antibiotic ointment if available
  \item Cover with sterile bandage
  \item Change daily or when dirty/wet
  \item Watch for signs of infection: increasing redness, warmth, swelling, pus, fever, red streaks extending from wound
\end{itemize}

\section{Review Questions}

\begin{enumerate}[leftmargin=*]
\item How do you classify a burn as superficial, partial-thickness, or full-thickness?
\item Why is the Rule of Nines useful, and what is it used for in adults?
\item When should a burn be cooled with running water, and for how long?
\item Why should you never attempt to remove clothing stuck to a burn?
\item What is the RICE protocol, and when is it appropriate?
\item A person falls and cannot move their legs. Why should they not be moved, and what should you do?
\end{enumerate}
\index{burns}
\index{thermal injury}
\index{Rule of Nines}
\index{burn extent}
\index{chemical burns}
\index{electrical burns}
\index{fracture management}
\index{open fracture}
\index{splinting}
\index{sprain}
\index{strain}
\index{RICE protocol}
\index{concussion}
\index{head injury}
\index{spinal injury}
\index{foreign body in eye}
\index{eye injury management}
\index{contusion}
\index{abrasion}
\index{open fracture}
\index{compound fracture}
\index{chest injury}
\index{sucking chest wound}
\index{occlusive dressing}
\index{pneumothorax}
\index{eye irrigation}